\section{Main Theorem}\label{sec:main}

\begin{theorem}[Conditional positive limiting loss]\label{thm:main}
Fix $\kappa\geq1$ and $q>0$.  Let
$r_0,C_{\mathrm{dim}},\delta,L_P,\zeta,C_T$ be positive quantities that
may depend only on $(\kappa,q)$, with $r_0\in\mathbb N$ and
\[
L_P<\delta/4,
\qquad
\zeta<\delta/4.
\]
Define
\[
\epsilon=
\left(\frac{\delta-L_P-\zeta}{\kappa^6C_T}\right)^2>0.
\]
For every $r\geq r_0$, every $n,k$ satisfying
Assumptions~\ref{assump:dimension} and \ref{assump:rank_window}, every
deterministic base triple satisfying
Assumption~\ref{assump:base_conditioning}, and the half-relaxed parallel ALS
trajectory under Assumptions~\ref{assump:gaussian_smoothing} and
\ref{assump:independent_initialization}, one has the event inclusion
\[
\mathsf C_2(\delta,L_P,\zeta,C_T)
\subseteq
\left\{
\begin{array}{l}
\displaystyle \lim_{t\to\infty}\mathcal L(X_t,Y_t,Z_t)
\text{ exists and is finite},\\[2mm]
\displaystyle \lim_{t\to\infty}\mathcal L(X_t,Y_t,Z_t)
\geq\epsilon\|T\|_F^2
\end{array}
\right\}.
\]
This is a deterministic implication on each realization in the displayed
event.  It makes no assertion that $\mathsf C_2$ is nonempty and no assertion
that it has positive probability.

There are no hidden constants in the displayed loss fraction.  Its exposed
dependence is exactly on $\kappa,\delta,L_P,\zeta,C_T$; $r,n,k$ enter only
through the theorem scope and certificate objects.  The probability mode is
event inclusion under the joint law, the residual floor used in the proof is
all-time, the conclusion is asymptotic, and all residuals and objectives use
the ambient Frobenius norm.

Separately from the preceding positive-Gaussian-smoothing event inclusion,
consider the deterministic algebraic zero-smoothing baseline in a tall
ambient space $n\geq r$: each of $\bar A,\bar B,\bar C$ has orthonormal
columns, every perturbation vector is set to zero, and hence
$T=\sum_{j=1}^r\bar a_j\otimes\bar b_j\otimes\bar c_j$.  In this baseline,
the ambient-to-coefficient map is generally rectangular, but it satisfies the
well-typed relations
\[
QT=D_r,
\qquad E_\rho=0,
\qquad \|Q\|_{\mathrm{op}}=1,
\qquad \|T\|_F=\|D_r\|_F.
\]
For any half-relaxed parallel ALS trajectory for this deterministic target
that satisfies the normalized entry-deficit, finite projector-path, and
finite represented-tensor-variation clauses, the same argument gives the
stronger baseline conclusion
\[
\lim_{t\to\infty}\mathcal L(X_t,Y_t,Z_t)
\geq(\delta-L_P)^2\|T\|_F^2.
\]
This last statement is an algebraic noiseless comparison, not an event under
Assumption~\ref{assump:gaussian_smoothing} or the positive-smoothing joint law
used in the first part of the theorem.
In particular, because the conditional implication holds for every positive
choice with the stated margins, it implies the existential constant statement
in the formalized goal, with all theorem constants depending only on
$(\kappa,q)$.
\end{theorem}

The unresolved source-level problem is to prove a uniform bound
$\mathbb P[\mathsf C_2(\delta,L_P,\zeta,C_T)]\geq p_0(\kappa,q)>0$ for
some admissible constants.  Theorem~\ref{thm:main} neither proves nor assumes
such a bound.
